% -------------------------------------------------
% PACKAGES SUPPLEMENTAIRES au besoin
% -------------------------------------------------

% Les paquets déjà chargés sont les suivants :
%-->    \usepackage[utf8]{inputenc} ou \usepackage[latin1]{inputenc}
%-->    \usepackage[T1]{fontenc}
%-->    \usepackage{lmodern}
%-->    \usepackage{enumerate}
%-->    \usepackage{amsfonts, amsmath, amssymb, stmaryrd, latexsym}
%-->    \usepackage{xspace, setspace}
%-->    \usepackage{array}
%-->    \usepackage[dvipsnames,usenames]{color}
%-->    \usepackage[table]{xcolor}
%-->    \usepackage{wrapfig, epsfig}
%-->    \usepackage[english,frenchb]{babel}
%-->    \usepackage{frbib} % bibliography en francais
%-->    \usepackage{fancyhdr}
%-->    \usepackage{geometry}
%-->    \usepackage{ulem} % for sout and underline
%-->    \usepackage{times,helvet,courier}
%-->    \usepackage{listings} (pour le code source)
%-->    \usepackage[pdftex, colorlinks, linkcolor=blue]{hyperref}
%-->    \usepackage[square, numbers]{natbib}
%-->    \usepackage[nomain, automake, abbreviations, nonumberlist, toc]{glossaries-extra} % abréviations
%-->    \usepackage{booktabs} 
%-->    \usepackage{doi}
%-->    \usepackage{refstyle}   % permet de faciliter les références croisées et le changement de langue
%-->    \usepackage[final]{pdfpages} % permet l'inclusion de pdf dans le document (voir contenu/articlePublie.tex)
%-->    \usepackage[a-1b]{pdfx}	

%\usepackage{biblatex}

%\usepackage{splitbib}
%\begin{category}[A]{First category}
%	\SBentries{deschenes98,BuDo.99-guide-B}
%\end{category}
%\begin{category}[B]{Second category}
%		\SBentries{Abr.96-BBook,Hoa.85-CSP,Jac.83-JSD,Mil.89-CCS}
%\end{category}


%% =====================================
%% Ajoutez, enlevez ou adaptez ICI
%% =====================================

%https://tug.org/FontCatalogue/paratypesansnarrow/
%\usepackage{PTSansNarrow} ptsn
%\newfontfamily\ptsansnarrow{PTSansNarrow-Regular.ttf}[
%    BoldFont = PTSansNarrow-Bold.ttf ,
%    Extension = .ttf
%]

\allowdisplaybreaks
\usepackage{xurl} %split url dans la biblio
\usepackage[nooneline]{subfigure}

%% pour l'ajout des abréviations
\usepackage[nomain, automake, abbreviations, nonumberlist, toc]{glossaries-extra}

\newglossary*{glossaire}{Glossaire}
\newglossary*{abreviationen}{AbreviationEn}
\makeglossaries

\usepackage{datetime}
%\usepackage{bigfoot} change the font
\usepackage[hang,bottom]{footmisc}
\usepackage{needspace}

\usepackage{paralist}
\usepackage{algorithm}

\usepackage{longtable}
\usepackage{threeparttablex}

%% permet de faire du pseudo code en francais - peut être modifié
\usepackage{algpseudocode}

\if@english 
	% defaut du paquet
\else % vous pouvez adapter au besoin
	\algrenewcommand\algorithmicwhile{\textbf{tant\_que}}
	\algnewcommand\algorithmiccase{\textbf{}}
	\algnewcommand\algorithmicswitch{\textbf{selon}}
	\algnewcommand\algorithmicelsif{\textbf{sinon si}}
	\algrenewcommand\algorithmicdo{}%\textbf{faire}
	\algrenewcommand\algorithmicif{\textbf{si}}
	\algrenewcommand\algorithmicelse{\textbf{sinon}}
	\algrenewcommand\algorithmicthen{\textbf{alors}}
	\algrenewcommand\algorithmicend{\textbf{fin}}
	\algrenewcommand\algorithmicfunction{\textbf{fonction}}
	\algrenewcommand\algorithmicfor{\textbf{pour}}
	\algrenewcommand\algorithmicrepeat{\textbf{répéter}}
	\algrenewcommand\algorithmicuntil{\textbf{jusqu'à}}
	\algrenewcommand\algorithmicreturn{\textbf{retourner}}
	\algrenewtext{EndFunction}{\algorithmicend}
	\algrenewtext{EndFor}{\algorithmicend\_\algorithmicfor}
	\algrenewtext{EndCase}{\algorithmicend\_\algorithmiccase}
	\algrenewtext{EndIf}{\algorithmicend\_\algorithmicif}
	\algrenewtext{EndSwitch}{\algorithmicend\_\algorithmicswitch}
	\algrenewtext{EndWhile}{\algorithmicend\_\algorithmicwhile}
	\algdef{SE}[WHILE]{While}{EndWhile}[1]{\algorithmicwhile\ #1\ \algorithmicdo}{\algorithmicend\_\algorithmicwhile}%
	\algdef{SE}[SWITCH]{Switch}{EndSwitch}[1]{\algorithmicswitch\ #1}{\algorithmicend\_\algorithmicswitch}%
	\algdef{SE}[CASE]{Case}{EndCase}[1]{\algorithmiccase\ #1 :}{\algorithmicend\_\algorithmiccase}%
	\algtext*{EndCase}
	\algrenewcommand\algorithmicrequire{\textbf{Pré-condition}}
	\algrenewcommand\algorithmicensure{\textbf{Post-condition}}
\fi

% -------------------------------------------------
% COMMANDES SUPPLEMENTAIRES
% -------------------------------------------------

%---- pour l'anglais ----
\newcommand{\ie}{{i.e.}\xspace}
\newcommand{\eg}{{e.g.}\xspace}
\newcommand{\cf}{{\it cf.}\xspace}

%---- pour le francais ----
\newcommand{\cad}{{c'est-\`a-dire}\xspace}
\newcommand{\Cad}{{C'est-\`a-dire}\xspace}
\newcommand{\etc}{{\it etc.}\xspace}

%---- Math ----
\newcommand{\N}{\ensuremath{\mathbb{N}}\xspace}
\newcommand{\Z}{\ensuremath{\mathbb{Z}}\xspace}

% Nouvelles d\'efintion d'environnement de th\'eor\`eme et de d\'efinition.
\newtheorem{frtheoreme}{Th\'eor\`eme}[section]
\newtheorem{frlemme}[frtheoreme]{Lemme}
\newtheorem{frprop}[frtheoreme]{Proposition}
\newtheorem{frcoro}[frtheoreme]{Corollaire}
\newtheorem{frdefinition}{D\'efinition}[section]

%---- Ajout autres memoires ----
\newcommand{\anglais}[1]{\textit{#1}}
\setabbreviationstyle[abbreviation]{long-short}

\newabbreviationstyle{longsansparenthese-court}
{%
  \renewcommand*{\CustomAbbreviationFields}{%
    name={\protect\glsxtrscfont\the\glsshorttok},
    sort={\the\glsshorttok},
    first={\the\glslongtok},
    firstplural={\the\glslongpltok},
    text={\protect\glsxtrscfont{\the\glsshorttok}},
    plural={\protect\glsxtrscfont{\the\glsshortpltok}},
    description={\the\glslongtok}}%
  \renewcommand*{\GlsXtrPostNewAbbreviation}{%
    \glssetattribute{\the\glslabeltok}{regular}{true}}%
}{}
\setabbreviationstyle[anglais]{longsansparenthese-court}
\newcommand{\glsentryshortwithlink}[1]{\glslink{#1}{\glsentryshort{#1}}}

%---- Ajout de Maxime ----

%\usepackage{natbib}
\usepackage{cancel} %\cancel{text}
\usepackage{amsthm}
\usepackage{amsmath}
\usepackage{amssymb}
\usepackage{cleveref}
\usepackage{etoolbox}
\usepackage{enumitem}
%\usepackage{caption}
\usepackage{adjustbox}
%\usepackage{lscape}
\usepackage{pdflscape}

%\usepackage{underscore} % permet d'utiliser _ in math mode.
\usepackage{threeparttable}
\renewcommand\operatorname[1]{\textit{#1}}

\renewcommand\labelitemi{{\boldmath$\cdot$}}

%\setcounter{tocdepth}{2} %Ne pas mettre plus de 2 de profond

\AtBeginEnvironment{align}{\setcounter{equation}{0}}
\counterwithout{equation}{chapter} % For report/book classes

\edef\hc{\string:} % \hc prints a normal colon
\setlength{\parindent}{0pt} %ligne 102

\newcommand{\HUGE}{\fontsize{40}{48}\selectfont}

\newcommand{\refenonce}[3]{%
\ifthenelse{\equal{#2}{}}%
{}%\textit{(paragraphe : #3)\hspace{-0.5cm}}}
{\leavevmode\\%
{\footnotesize\fontfamily{ptm}\selectfont(voir énoncé[s] lié[s] : #1 p.\href{#1}{\pageref{enonce:#2}})\hspace{-0.52cm}}}%\hspace{-0.52cm}}}%
}

\usepackage[defaultlines=2,all]{nowidow}

\usepackage{expl3}

\newcounter{paragraphenumero}

\makeatletter
\renewcommand{\theparagraphenumero}{\two@digits{\value{paragraphenumero}}}
\makeatother


\newcommand{\paragraphe}[4]{%
%\interfootnotelinepenalty=10000% ne règle pas tout à fait le saut de page entre le numéro et le reste du texte.
\clubpenalty=10000%   % Prohibits a single orphan line at the bottom of the page
\widowpenalty=10000%  % Prohibits a single widow line at the top of the next page
\ifthenelse{\equal{#1}{}}%
{}%
{\ifthenelse{\equal{#3}{}}%
{}%
{\stepcounter{paragraphenumero}\hspace*{-0.95cm}%1.33
{\fontfamily{ptm}\selectfont{P\theparagraphenumero.\phantom{a}}}}}%
#2%
\refenonce{#3}{#4}{#1}%
\label{paragraphe:#1}%
\vspace{\baselineskip}\par
}%
%\newline
%\vspace{0.5\baselineskip}\par
%\end{samepage}

\newcommand{\enonce}[5]{}
\newcommand{\enoncewithoutexp}[5]{}
\newcommand{\explication}[5]{}

%\newcommand{\sep}[0]{\phantom{a}\newline}
%\newcommand{\sep}[0]{\smallskip}
\newcommand{\sep}[0]{\hfill\par}

\usepackage{chngcntr} % Éviter que les note de bas de pages recommencent,
\counterwithout{footnote}{chapter}

%• Times New Roman, 12 points • Times, 12 points • Helvetica, 11 points • Geneva, 10 points
%https://www.overleaf.com/learn/latex/Font_typefaces
\newcommand{\fl}[1]{%
{\fontfamily{ptm}\selectfont#1}%ptm
}

\newcommand{\enfr}[2]{
[Traduction libre] «~#2~» \footnote{\textit{«~#1~»}}%
}

\newcommand{\enfrtest}[2]{
[Traduction libre] #2
\footnote{«~#1~»}
}

\newcommand{\microskip}{\vspace{1.5pt plus 0.5pt minus 0.5pt}}

\newcommand{\quotefoot}[1]{%
\begin{quote}%
#1%
\end{quote}%
}

\newcommand{\listfoot}[1]{%
\textit{%
\begin{itemize}[nosep,topsep=\smallskipamount]%
#1%
\end{itemize}%
}%
}

\newcommand{\quoteenfr}[3]{%
\begin{quote}%
[Traduction libre] #2 \footnote{%
%\clubpenalty=10000%   % Prohibits a single orphan line at the bottom of the page
%\widowpenalty=10000%  % Prohibits a single widow line at the top of the next page
%\interfootnotelinepenalty=0
%\clubpenalty=0%   % Prohibits a single orphan line at the bottom of the page
%\widowpenalty=0%  % Prohibits a single widow line at the top of the next page
\quotefoot{\textit{#1}}} #3%
\end{quote}%
}

%Mettre l'italic comme ceci parce que l'affichage de la première bullet n'apparaitra pas.
\newcommand{\quoteenfrlist}[3]{
\begin{quote}
[Traduction libre]
\begin{itemize}[leftmargin=*, labelindent=-0.62cm]
#2
\footnote{%
\listfoot{%
#1%
}
}
{\normalfont #3}
\end{itemize}%
\end{quote}%
}

\newcommand{\quoteenfrlistcustom}[4]{
\begin{quote}
#4
\begin{itemize}[leftmargin=*, labelindent=-0.62cm]
#2
\footnote{%
\textit{%
\begin{itemize}
#1
\end{itemize}
}
}
{\normalfont #3}
\end{itemize}
\end{quote}

}

% Provenance de https://tex.stackexchange.com/questions/3389/how-to-indicate-elision-in-a-quotation
%\newcommand*\elide{\textup{[\,\ldots\,]}\xspace}
\newcommand*\elide{\textup{[...]}\xspace}


\newcommand{\quotefr}[2]{
\begin{quote}
#1
#2
\end{quote}
}

%#1 : clé, #2 : le terme, #3 : la définition, #4 : le séparateur 
\newcommand{\glo}[4]{
#2#4#3
\newabbreviation[type=glossaire]{#1}{#2}{#3}
}

%#1 : clé, #2 : l'acronyme, #3 : le terme
\newcommand{\acr}[3]{
#2 #3
\newabbreviation[type=acronymtype]{#1}{#2}{#3}
}

% Normalement le template devrait déjà gérer les quotes selon les normes du département.
% Il doit y avoir une erreur dans les normes du département.
% Citations dans le texte
% Les citations de plus de deux lignes sont écrites entièrement en retrait de 2,5 cm de la
% marge de gauche, à interligne simple, placées entre guillemets, et précédées et suivies % d’un saut de ligne.
\newenvironment{quo}[1]
{
\list{}{\leftmargin=2.5cm \rightmargin=0pt} % On définit le retrait de 2,5cm
  \item\relax
  \begin{spacing}{1}%
}
{
\end{spacing}\endlist
}

% titre, lien, 
%\fig
%{\textit{Quality attributes mentioned on each publication.}}
%{figure/chap1/documentation/santos.png}
%{Source~: tableau 1 \cite{santos_review_2020}}
%{}

\newcommand{\finfigtab}[2]{
\if\relax\detokenize{#1}\relax
  % vide
\else
  \break\onehalfspacing\raggedright
\fi
\if\relax\detokenize{#2}\relax
  % vide
\else
  Traduction libre\\
\fi
\if\relax\detokenize{#1}\relax
  % vide
\else
  #1
\fi
}

\newcommand{\tab}[6]{
\begin{table}[#6]
\centering
\begin{threeparttable}
\caption[#1]{#1 #3 #4}
\begin{tabular}{c}
\includegraphics[width=0.99\linewidth]{#2}
\end{tabular}
#5
\end{threeparttable}
%\finfigtab{#3}{#4}
\end{table}
}

\newcommand{\tabmedium}[6]{
\begin{table}[#6]
\caption[#1]{#1 #3 #4}
\centering
\begin{threeparttable}
\begin{minipage}{0.8\linewidth}
\includegraphics[width=1\linewidth]{#2}
\end{minipage}
#5
\end{threeparttable}
\end{table}
}

\newcommand{\taben}[7]{
\begin{table}[#6]
\caption[#7]{#7 #3 #4}
\centering
\begin{threeparttable}
\fbox{%
\begin{minipage}{0.975\linewidth}
\centering
{\normalsize{\selectfont #1}}
\vspace{0.5em}
\includegraphics[width=1\linewidth]{#2}
\end{minipage}
}
#5
\end{threeparttable}
\end{table}
}

\newcommand{\tabmediumen}[7]{
\begin{table}[#6]
\caption[#7]{#7 #3 #4}
\centering
\begin{threeparttable}
\fbox{%
\begin{minipage}{0.8\linewidth}
\centering
{\normalsize{\selectfont #1}}
\vspace{0.5em}
\includegraphics[width=1\linewidth]{#2}
\end{minipage}
}
#5
\end{threeparttable}
\end{table}
}

\newcommand{\tabsmallen}[7]{
\begin{table}[#6]
\caption[#7]{#7 #3 #4}
\centering
\begin{threeparttable}
\fbox{%
\begin{minipage}{0.6\linewidth}
\centering
{\normalsize{\selectfont #1}}
\vspace{0.5em}
\includegraphics[width=1\linewidth]{#2}
\end{minipage}
}
#5
\end{threeparttable}
\end{table}
}

\newcommand{\tabsmallens}[7]{
\begin{table}[#6]
\caption[#7]{#7 #3 #4}
\centering
\begin{threeparttable}
\fbox{%
\begin{minipage}{0.6\linewidth}
\includegraphics[width=1\linewidth]{#2}
\end{minipage}
}
#5
\end{threeparttable}
\end{table}
}

\newcommand{\tabmen}[7]{
\begin{table}[#6]
\caption[#7]{#7 #3 #4}
\centering
\begin{threeparttable}
\fbox{%
\begin{minipage}{0.90\linewidth}
\centering
{\normalsize{\selectfont #1}}
\vspace{0.5em}
\includegraphics[width=0.95\linewidth]{#2}
\end{minipage}
}
#5
\end{threeparttable}
\end{table}
}

\newcommand{\tabtab}[6]{
\begin{table}[#6]
\centering
\begin{threeparttable}
\caption{#1}
#2
#5
\end{threeparttable}
%\finfigtab{#3}{#4}
\end{table}
}

\newcommand{\tabtaben}[7]{
\begin{table}[#6]
\caption[#7]{#7 #3 #4}
\centering
\begin{threeparttable}
\fbox{\adjustbox{width=0.97\linewidth}{%
#2%
}%
}%
#5
\end{threeparttable}
\end{table}
}

\newcommand{\tabtabensmall}[7]{
\begin{table}[#6]
\caption[#7]{#7 #3 #4}
\centering
\begin{threeparttable}
\fbox{\adjustbox{width=0.80\linewidth}{%
#2%
}%
}%
#5
\end{threeparttable}
\end{table}
}

\newcommand{\fig}[6]{
\begin{figure}[#6]
\centering
\begin{threeparttable}
\begin{tabular}{c}
\includegraphics[width=0.985\linewidth]{#2}
\end{tabular}
\captionsetup{width=\linewidth}
\caption[#1]{#1 #3 #4}
#5
\end{threeparttable}
\end{figure}
}

\newcommand{\figen}[7]{
\begin{figure}[#6]
\centering
\begin{threeparttable}
\fbox{%
\begin{tabular}{c}%
\adjustbox{width=0.95\linewidth}{%
\includegraphics{#2}%
}\\[1ex]%
{\parbox{0.955\linewidth}{\centering #1}}%
\end{tabular}%
}%
#5%
\caption[#7]{#7 #3 #4}%
\end{threeparttable}
\end{figure}
}


%ajustbox doit être avant caption dans threeparttable puisqu'il y a un calcul automatique de la taille qui est fait dans cette lib. Donc, il est impossible d'utiliser la même stratégie pour un tabsmall.
\newcommand{\figsmall}[6]{
\begin{figure}[#6]
\centering
\begin{threeparttable}
    \begin{adjustbox}{width=0.5\linewidth}
        \begin{tabular}{c}
        \includegraphics[]{#2}
        \end{tabular}
    \end{adjustbox}
    \caption[#1]{#1 #3 #4}
    #5
\end{threeparttable}
%\finfigtab{#3}{#4}
\end{figure}
}

\newcommand{\figms}[6]{
\begin{figure}[#6]
\centering
\begin{threeparttable}
    \begin{adjustbox}{width=0.60\linewidth}
        \begin{tabular}{c}
        \includegraphics[]{#2}
        \end{tabular}
    \end{adjustbox}
    \caption[#1]{#1 #3 #4}
    #5
\end{threeparttable}
%\finfigtab{#3}{#4}
\end{figure}
}

\newcommand{\figsmallen}[7]{
\begin{figure}[#6]
\centering
\begin{threeparttable}
    \fbox{%
    \begin{tabular}{c}
      \adjustbox{width=0.5\linewidth}{%
        \includegraphics{#2}%
      }\\[1ex]
      {\parbox{0.955\linewidth}{\centering #1}}
    \end{tabular}
    }
    #5
    \caption[#7]{#7 #3 #4}
\end{threeparttable}
\end{figure}
}

\newcommand{\figmsen}[7]{
\begin{figure}[#6]
\centering
\begin{threeparttable}
    \fbox{%
    \begin{tabular}{c}
      \adjustbox{width=0.65\linewidth}{%
        \includegraphics{#2}%
      }\\[1ex]
      {\parbox{0.955\linewidth}{\centering #1}}
    \end{tabular}
    }
    #5
    \caption[#7]{#7 #3 #4}
\end{threeparttable}
\end{figure}
}

\newcommand{\figggmen}[7]{
\begin{figure}[#6]
\centering
\begin{threeparttable}
    \fbox{%
    \begin{tabular}{c}
      \adjustbox{width=0.95\linewidth}{%
        \includegraphics{#2}%
      }\\[1ex]
      {\parbox{0.955\linewidth}{\centering #1}}
    \end{tabular}
    }
    #5
    \caption[#7]{#7 #3 #4}
\end{threeparttable}
\end{figure}
}

\newcommand{\figgmen}[7]{
\begin{figure}[#6]
\centering
\begin{threeparttable}
    \fbox{%
    \begin{tabular}{c}
      \adjustbox{width=0.85\linewidth}{%
        \includegraphics{#2}%
      }\\[1ex]
      {\parbox{0.955\linewidth}{\centering #1}}
    \end{tabular}
    }
    #5
    \caption[#7]{#7 #3 #4}
\end{threeparttable}
\end{figure}
}

\newcommand{\figgmmen}[7]{
\begin{figure}[#6]
\centering
\begin{threeparttable}
    \fbox{%
    \begin{tabular}{c}
      \adjustbox{width=0.75\linewidth}{%
        \includegraphics{#2}%
      }\\[1ex]
      {\parbox{0.955\linewidth}{\centering #1}}
    \end{tabular}
    }
    #5
    \caption[#7]{#7 #3 #4}
\end{threeparttable}
\end{figure}
}

\newcommand{\figmediumen}[7]{
\begin{figure}[#6]
\centering
\begin{threeparttable}
    \fbox{%
    \begin{tabular}{c}
      \adjustbox{width=0.75\linewidth}{%
        \includegraphics{#2}%
      }\\[1ex]
      {\parbox{0.955\linewidth}{\centering #1}}
    \end{tabular}
    }
    #5
    \caption[#7]{#7 #3 #4}
\end{threeparttable}
\end{figure}
}

\newcommand{\figens}[7]{
\begin{figure}[#6]
\centering
\begin{threeparttable}
    \fbox{%
    \begin{tabular}{c}
      \adjustbox{width=0.96\linewidth}{%
        \includegraphics{#2}%
      }
    \end{tabular}
    }
    #5
    \caption[#7]{#7 #3 #4}
\end{threeparttable}
\end{figure}
}

\newcommand{\figmediumens}[7]{
\begin{figure}[#6]
\centering
\begin{threeparttable}
    \fbox{%
    \begin{tabular}{c}
      \adjustbox{width=0.80\linewidth}{%
        \includegraphics{#2}%
      }
    \end{tabular}
    }
    #5
    \caption[#7]{#7 #3 #4}
\end{threeparttable}
\end{figure}
}

\newcommand{\figmedium}[6]{
\begin{figure}[#6]
\centering
\begin{threeparttable}
    \begin{adjustbox}{width=0.75\linewidth}
        \begin{tabular}{c}
        \includegraphics[]{#2}
        \end{tabular}
    \end{adjustbox}
    \caption[#1]{#1 #3 #4}
    #5
\end{threeparttable}
%\finfigtab{#3}{#4}
\end{figure}
}

\newcommand{\figgm}[6]{
\begin{figure}[#6]
\centering
\begin{threeparttable}
    \begin{adjustbox}{width=0.90\linewidth}
        \begin{tabular}{c}
        \includegraphics[]{#2}
        \end{tabular}
    \end{adjustbox}
    \caption[#1]{#1 #3 #4}
    #5
\end{threeparttable}
%\finfigtab{#3}{#4}
\end{figure}
}

%\usepackage[paper=A4]{typearea} réduit l'entête

% Provient de https://tex.stackexchange.com/questions/9071/how-to-translate-and-rotate-the-heading-of-landscaped-pages
\usepackage{geometry}

\usepackage{fancyhdr}
\pagestyle{fancyplain}

%\newlength{\tempwidth} % Defines the temporary variable
%\setlength{\tempwidth}{\textwidth}
%\setlength{\textwidth}{\headwidth}
%\setlength{\headwidth}{\tempwidth}
  
\fancypagestyle{landscapestyle}{
  \fancyhead{}
  \fancyfoot{}
}

%\makebox[0pt][l]{\hspace*{1cm}\vspace*{1cm}\rotatebox{90}{Left header in margin}} 
%\makebox[0pt][r]{\hspace*{3cm}\rotatebox{90}{\thepage}}%

\usepackage[absolute,overlay]{textpos}

%Dans les normes du départment, cela prend au moins 3cm dans la marge de gauche pour la reliure, puisque c'est en mode paysage la reliure devrait être en haut.
%left=3cm,right=2.5cm,top=3.75cm,bottom=2.5cm,
%29.7 × 21 cm paysage A4
\newcommand{\tablandscape}[6]{
{%
\newgeometry{left=3cm,right=2.5cm,top=5cm,bottom=2.5cm}
\begin{landscape}
\pagestyle{empty}
\begin{table}[#6]
\begin{textblock*}{1cm}(1cm,14cm)
\makebox[0pt][l]{\hspace*{1cm}\vspace*{1cm}\rotatebox{90}{\rightmark\hfill}}
\end{textblock*}
\begin{threeparttable}
\caption{#1}
\begin{tabular}{c}
\includegraphics[width=1.4\textwidth]{#2}
\end{tabular}
\setstretch{0.8}
#5
\end{threeparttable}
\centering
\begin{textblock*}{1cm}(17.5cm,13.7cm)
\makebox[0pt][l]{\hspace*{1cm}\vspace*{1cm}\rotatebox{90}{\thepage}}
\end{textblock*}
\end{table}
\end{landscape}
}%
}

\newcommand{\tablandscapec}[6]{
{%
\newgeometry{left=3cm,right=2.5cm,top=5cm,bottom=2.5cm}
\begin{landscape}
\pagestyle{empty}
\begin{table}[]
\begin{textblock*}{1cm}(1cm,14cm)
\makebox[0pt][l]{\hspace*{1cm}\vspace*{1cm}\rotatebox{90}{\rightmark\hfill}}
\end{textblock*}
\begin{threeparttable}
\addtocounter{table}{-1}
\caption[]{#1}
\begin{tabular}{c}
\includegraphics[width=1.4\textwidth]{#2}
\end{tabular}
\setstretch{0.8}
#5
\end{threeparttable}
\centering
\begin{textblock*}{1cm}(17.5cm,13.7cm)
\makebox[0pt][l]{\hspace*{1cm}\vspace*{1cm}\rotatebox{90}{\thepage}}
\end{textblock*}
\end{table}
\end{landscape}
}%
}

\newcommand{\tabletwo}[7]{
{%
\newgeometry{left=3cm,right=2.5cm,top=3.75cm,bottom=2.5cm}
\begin{landscape}
\pagestyle{empty}
\begin{longtable}{c}
\begin{textblock*}{1cm}(1cm,14cm)
\makebox[0pt][l]{\hspace*{1cm}\vspace*{1cm}\rotatebox{90}{\rightmark\hfill}}

\begin{threeparttable}
\caption{#1}

\caption{Two images across two pages} \\

% First page
\includegraphics[width=\textwidth]{image1} \\
(a) First image \\

\endfirsthead

\multicolumn{1}{c}{\textit{Table \thetable\ continued}} \\

% Second page
\includegraphics[width=\textwidth]{image2} \\
(b) Second image \\

\end{longtable}

\begin{table}[#6]

\end{textblock*}
\begin{threeparttable}
\caption{#1}
\begin{tabular}{c}
\includegraphics[width=1.4\textwidth]{#2}
\end{tabular}
\setstretch{0.8}
#5
\end{threeparttable}
\finfigtab{#3}{#4}
\centering
\begin{textblock*}{1cm}(17.5cm,13.7cm)
\makebox[0pt][l]{\hspace*{1cm}\vspace*{1cm}\rotatebox{90}{\thepage}}
\end{textblock*}
\end{table}
\end{landscape}
}%
}

\newcommand{\tablandscapeextra}[6]{
{%
\newgeometry{left=1.7cm,right=1cm,top=0.5cm,bottom=0.5cm}
\begin{landscape}
\begin{table}[#6]
\begin{textblock*}{1cm}(0cm,13cm)
\makebox[0pt][l]{\hspace*{1cm}\vspace*{1cm}\rotatebox{90}{\rightmark\hfill}}
\end{textblock*}
\begin{threeparttable}
\caption{#1}
\begin{tabular}{c}
\includegraphics[width=1.4\textwidth]{#2}
\end{tabular}
\setstretch{0.6}
#5
\end{threeparttable}
\centering
\begin{textblock*}{1cm}(17.5cm,13.7cm)
\makebox[0pt][l]{\hspace*{2cm}\vspace*{2cm}\rotatebox{90}{\thepage}}
\end{textblock*}
\end{table}
\end{landscape}
\clearpage 
\restoregeometry 
\clearpage % Force an immediate page break to lock in the restored geometry
\pagestyle{fancyplain}
\thispagestyle{fancyplain}
}%
}

\newcommand{\figlandscape}[6]{
{%
\newgeometry{left=3cm,right=2.5cm,top=5cm,bottom=2.5cm}
\begin{landscape}
\pagestyle{empty}
\begin{figure}[#6]
\begin{textblock*}{1cm}(1cm,16cm)
\makebox[0pt][l]{\hspace*{1cm}\vspace*{1cm}\rotatebox{90}{\rightmark\hfill}}
\end{textblock*}
\begin{threeparttable}
\begin{tabular}{c}
\includegraphics[width=1.4\textwidth]{#2}
\end{tabular}
#5
\caption{#1}
\end{threeparttable}
\finfigtab{#3}{#4}
\centering
\begin{textblock*}{1cm}(17.5cm,13.7cm)
\makebox[0pt][l]{\hspace*{1cm}\vspace*{1cm}\rotatebox{90}{\thepage}}
\end{textblock*}
\end{figure}
\end{landscape}
\restoregeometry
\pagestyle{fancyplain}
}%
}

\newcommand{\figlandscapel}[6]{
{%
\newgeometry{left=3cm,right=2.5cm,top=5cm,bottom=2.5cm}
\begin{landscape}
\pagestyle{empty}
\begin{figure}[#6]
\begin{textblock*}{1cm}(1cm,10cm)
\makebox[0pt][l]{\hspace*{1cm}\vspace*{1cm}\rotatebox{90}{\rightmark\hfill}}
\end{textblock*}
\begin{threeparttable}
\begin{tabular}{c}
\includegraphics[width=1.2\textwidth]{#2}
\end{tabular}
#5
\caption{#1}
\end{threeparttable}
\centering
\begin{textblock*}{1cm}(17.5cm,13.7cm)
\makebox[0pt][l]{\hspace*{1cm}\vspace*{1cm}\rotatebox{90}{\thepage}}
\end{textblock*}
\end{figure}
\end{landscape}
\restoregeometry
\pagestyle{fancyplain}
}%
}

\newcommand{\figlandscapeextra}[6]{
{%
\newgeometry{left=1.2cm,right=1.5cm,top=0.5cm,bottom=0.5cm}
\begin{landscape}
%\pagestyle{empty}
\begin{figure}[#6]
\begin{textblock*}{1cm}(1cm,14cm)
\makebox[0pt][l]{\hspace*{1cm}\vspace*{1cm}\rotatebox{90}{\rightmark\hfill}}
\end{textblock*}
\begin{threeparttable}
\caption{#1}
\begin{tabular}{c}
\includegraphics[width=1.4\textwidth]{#2}
\end{tabular}
#5
\end{threeparttable}
\finfigtab{#3}{#4}
\centering
\begin{textblock*}{1cm}(17.5cm,13.7cm)
\makebox[0pt][l]{\hspace*{1cm}\vspace*{1cm}\rotatebox{90}{\thepage}}
\end{textblock*}
\end{figure}
\end{landscape}
\restoregeometry
%\pagestyle{fancyplain}
}%

}

\newcommand{\titleen}[2]{%
\textit{#1}%\footnote{[Traduction libre] #2}
}

\newcommand{\titlefr}[1]{%
%{\fontfamily{qpl}\fontseries{sb}\selectfont#1}%
\textit{#1}%
}

\newcommand{\emphase}[1]{%
{\fontfamily{bch}\fontsize{11}{12}\selectfont#1}%
}

\newcommand{\titlefig}[1]{%
{\fontfamily{PTSansNarrow-TLF}\selectfont#1}%
}

%anglais, acronyme, francais
\newcommand{\delabbrev}[3]{%
\textit{#1}%
}

\newcommand{\acroIn}[2]{
\ifthenelse{\equal{#2}{np}}%
{#1}
{}
\ifthenelse{\equal{#2}{nca}}%
{\textit{#1}}
{}
\ifthenelse{\equal{#2}{nc}}%
{#1}
{}
}

%acronyme 
% #3 : nompropre (np), nomcommunanglais (nca), nomcommun (nc)
% #4 : complet (c), acro (a), nom(n)
\newcommand{\acro}[4]{%
\ifthenelse{\equal{#4}{c}}%
{\acroIn{#1}{#3}(#2)}
{}
\ifthenelse{\equal{#4}{n}}%
{\acroIn{#1}{#3}}
{}
\ifthenelse{\equal{#4}{a}}%
{#2}
{}
}

\newcommand{\figannexe}[1]{}

%limite de 9 arguments en lateX
\NewDocumentCommand{\ccite}{m O{} O{} O{} O{} O{} O{} O{} O{}}{%
\ifthenelse{\equal{#9}{}}%
{%
\ifthenelse{\equal{#7}{}}%
{%
\ifthenelse{\equal{#5}{}}%
{%
\ifthenelse{\equal{#3}{}}%
{%
\ifthenelse{\equal{#1}{}}%
{}%
{\cite{#1}%
}%
}%
{\cite{#1, #3}}%
}%
{\cite{#1, #3, #5}}%
}%
{\cite{#1, #3, #5, #7}}%
}%
{\cite{#1, #3, #5, #7, #9}}%
}%


%Pour conserver un vocabulaire traduit.
\newcommand{\sfren}[2]{
#1%
}

%\paragraphe{}{%
%\begin{textblock*}{1cm}(10.6cm,25.3cm)%
%\thepage%
%\end{textblock*}%
%\vspace{-1\baselineskip}
%}{}{}